%!TEX root = ../../main.tex
%% ===============================================================
%% 5.2 정의
%% 파일: chapters/05_label/02_definition.tex
%% 상위: chapters/05_label/chapter.tex
%% 정의 라벨: sec:label-definition, def:label, eq:exchange
%% 참조 라벨: eq:exchange, tab:label-weights
%% 포함 객체: tables/tab_label_weights
%% 인용: bahdanau2015neural
%% 상태: 초안 (docs/OUTLINE.md 참고)
%% ===============================================================

\section{정의}
\label{sec:label-definition}

\begin{definition}[교환가치 라벨]
\label{def:label}
라벨 창 $W_{\mathrm{lab}}$ 안의 사건 집합을 $U$라 하자. 각 사건 $u \in U$는 부호 $\sigma_u \in \{+1, -1\}$(블루 팀에 유리하면 $+1$), 가치 $v_u \ge 0$, 우선순위 $p_u \ge 0$을 갖는다. 주의 가중치 $\alpha_u = \softmax_u(\beta p_u) = e^{\beta p_u}/\sum_{u'\in U} e^{\beta p_{u'}}$로 정의하고 교환가치 점수를
\begin{equation}
  S(U) = \sum_{u \in U} \alpha_u\, \sigma_u\, v_u
  \label{eq:exchange}
\end{equation}
라 하면, 라벨은 $y = \ind[S(U) > 0]$이다. $\beta = 2.0$이다.
\end{definition}

식 \eqref{eq:exchange}의 softmax 가중은 Bahdanau 등 \cite{bahdanau2015neural}의 주의 메커니즘과 같은 형태로, 우선순위가 높은 사건(셧다운, 오브젝트)이 점수를 지배하도록 한다. $\beta \to 0$이면 균등 가중(단순 가치 합)에, $\beta \to \infty$이면 최고 우선순위 사건 하나가 결과를 결정하는 극한에 수렴한다. 사건의 가치 $v_u$는 표 \ref{tab:label-weights}의 도메인 가중치를 사건의 속성(킬 여부, 셧다운 여부, 연속 킬, 어시스트 수, 현상금, 오브젝트 등급, 라인 우선권, 특수 보너스)에 곱해 합산한 복합값이다.

%% ---------------------------------------------------------------
%% table 객체: tab:label-weights
%% 파일: tables/tab_label_weights.tex
%% 사용처: chapters/05_label/02_definition.tex  (%% ---------------------------------------------------------------
%% table 객체: tab:label-weights
%% 파일: tables/tab_label_weights.tex
%% 사용처: chapters/05_label/02_definition.tex  (%% ---------------------------------------------------------------
%% table 객체: tab:label-weights
%% 파일: tables/tab_label_weights.tex
%% 사용처: chapters/05_label/02_definition.tex  (\input{tables/tab_label_weights})
%% 규칙: 캡션·라벨·내용을 이 파일 안에서만 관리한다. 수치는 config/numbers.tex 매크로를 쓴다.
%% ---------------------------------------------------------------
\begin{table}[htbp]
  \centering
  \caption{교환가치 라벨의 사건 가중치}
  \label{tab:label-weights}
  \small
  \begin{tabular}{@{}llr@{}}
    \toprule
    속성 & 설명 & 가중치 \\
    \midrule
    kill & 챔피언 킬 & 1.00 \\
    shutdown & 셧다운 킬(연승 중인 상대 처치) & 1.60 \\
    streak & 연속 킬 & 0.35 \\
    assist & 어시스트(1 건당) & 0.20 \\
    bounty & 현상금 골드 & 0.30 \\
    objective & 오브젝트(용·바론·전령·포탑 등, 등급별) & 1.10 \\
    lane & 라인 관련 사건(포탑 방패 등) & 0.25 \\
    \midrule
    $\beta$ & 주의 온도 & 2.0 \\
    \bottomrule
  \end{tabular}
\end{table}
)
%% 규칙: 캡션·라벨·내용을 이 파일 안에서만 관리한다. 수치는 config/numbers.tex 매크로를 쓴다.
%% ---------------------------------------------------------------
\begin{table}[htbp]
  \centering
  \caption{교환가치 라벨의 사건 가중치}
  \label{tab:label-weights}
  \small
  \begin{tabular}{@{}llr@{}}
    \toprule
    속성 & 설명 & 가중치 \\
    \midrule
    kill & 챔피언 킬 & 1.00 \\
    shutdown & 셧다운 킬(연승 중인 상대 처치) & 1.60 \\
    streak & 연속 킬 & 0.35 \\
    assist & 어시스트(1 건당) & 0.20 \\
    bounty & 현상금 골드 & 0.30 \\
    objective & 오브젝트(용·바론·전령·포탑 등, 등급별) & 1.10 \\
    lane & 라인 관련 사건(포탑 방패 등) & 0.25 \\
    \midrule
    $\beta$ & 주의 온도 & 2.0 \\
    \bottomrule
  \end{tabular}
\end{table}
)
%% 규칙: 캡션·라벨·내용을 이 파일 안에서만 관리한다. 수치는 config/numbers.tex 매크로를 쓴다.
%% ---------------------------------------------------------------
\begin{table}[htbp]
  \centering
  \caption{교환가치 라벨의 사건 가중치}
  \label{tab:label-weights}
  \small
  \begin{tabular}{@{}llr@{}}
    \toprule
    속성 & 설명 & 가중치 \\
    \midrule
    kill & 챔피언 킬 & 1.00 \\
    shutdown & 셧다운 킬(연승 중인 상대 처치) & 1.60 \\
    streak & 연속 킬 & 0.35 \\
    assist & 어시스트(1 건당) & 0.20 \\
    bounty & 현상금 골드 & 0.30 \\
    objective & 오브젝트(용·바론·전령·포탑 등, 등급별) & 1.10 \\
    lane & 라인 관련 사건(포탑 방패 등) & 0.25 \\
    \midrule
    $\beta$ & 주의 온도 & 2.0 \\
    \bottomrule
  \end{tabular}
\end{table}

